%==============================================================================
%  Chapter 3 — Methodology
%
%  Structure follows the standard qualitative-research reporting template:
%   1. Research philosophy / paradigm
%   2. Research design (comparative qualitative case study)
%   3. Case selection and units of analysis
%   4. Data sources and collection (interviews, documents, observation)
%   5. Data analysis (coding, themes, framework application)
%   6. Trustworthiness (credibility, transferability, dependability,
%      confirmability)
%   7. Ethical considerations
%   8. Reflexivity and limitations
%==============================================================================
\chapter{Methodology}
\label{ch:methodology}

\section*{Chapter overview}
\todoinline{One paragraph: the chapter justifies a comparative qualitative
case study and details the data and analytical procedures.}

%==============================================================================
\section{Research philosophy}
\label{sec:meth-philosophy}
\todoinline{State your ontological and epistemological position
(interpretivist / constructionist is the typical pairing for an
institutional-logics study). Justify in 1--2 paragraphs.}

%==============================================================================
\section{Research design}
\label{sec:meth-design}

\subsection{Comparative qualitative case study}
\label{sec:meth-design-comparative}
\todoinline{Explain why a comparative case design is appropriate
(\textcite{eisenhardt1989building}; \textcite{yin2018case}). Distinguish your
\emph{cases} (Pathfinder, STEP) from your \emph{units of analysis}
(individual funding decisions, programme manager practices, etc.).}

\subsection{Case selection rationale}
\label{sec:meth-cases}
\todoinline{Use a polar / contrast logic if possible: same institutional
home, contrasting investment logics. Justify the comparability of the two
cases (timing, governance, instrument type).}

%==============================================================================
\section{Data sources and collection}
\label{sec:meth-data}

\subsection{Document corpus}
\label{sec:meth-docs}
\todoinline{List document categories (EIC Work Programmes 2021--2025, STEP
Regulation, evaluation guides, EIC Board reports, EU Court of Auditors
reports, news coverage). Show how you compiled and version-controlled the
corpus.}

\subsection{Semi-structured interviews}
\label{sec:meth-interviews}
\todoinline{Number of interviews, sampling strategy (purposive +
snowball), informant categories (programme managers, evaluators,
applicants, ecosystem intermediaries), interview duration, language,
recording and transcription practice.}

\begin{table}[H]
\centering
\caption[Interview informant overview]{Overview of interview informants.}
\label{tab:informants}
\begin{tabularx}{\textwidth}{@{} l l X r @{}}
\toprule
\textbf{ID} & \textbf{Role category} & \textbf{Affiliation type} & \textbf{Duration (min)} \\
\midrule
P1  & Programme Manager   & EISMEA          & --- \\
P2  & Evaluator           & Independent     & --- \\
P3  & Applicant (funded)  & University spin-out & --- \\
% Add rows as interviews are completed.
\bottomrule
\end{tabularx}
\end{table}

\subsection{Secondary data}
\label{sec:meth-secondary}
\todoinline{Funded-project databases (CORDIS, EIC dashboard), news media,
LinkedIn announcements, etc.}

%==============================================================================
\section{Data analysis}
\label{sec:meth-analysis}

\subsection{Coding strategy}
\label{sec:meth-coding}
\todoinline{Hybrid deductive--inductive coding using NVivo / Atlas.ti /
manual. First-order codes from interviewees' language; second-order themes
inferred; aggregate dimensions tied to your analytical framework
(Gioia methodology — \textcite{gioia2013seeking}).}

\subsection{Within-case and cross-case analysis}
\label{sec:meth-cross-case}
\todoinline{Describe the two analytical passes: (i) within-case narrative
for each instrument; (ii) structured comparison along the dimensions of the
framework.}

%==============================================================================
\section{Trustworthiness and quality criteria}
\label{sec:meth-quality}
\todoinline{Address Lincoln \& Guba's four criteria: credibility (member
checks, triangulation), transferability (thick description),
dependability (audit trail), confirmability (reflexive memos).}

%==============================================================================
\section{Ethical considerations}
\label{sec:meth-ethics}
\todoinline{Informed consent, anonymisation, GDPR compliance, data storage
policy, sensitivity around speaking on behalf of an EU institution.}

%==============================================================================
\section{Reflexivity}
\label{sec:meth-reflexivity}
\todoinline{Reflect on your own positionality — DTU affiliation, prior
exposure to EIC processes, etc. — and how it shaped access and
interpretation.}

%==============================================================================
\section{Methodological limitations}
\label{sec:meth-limitations}
\todoinline{Anticipate the limitations the reader will raise — small N,
contemporaneous data on a moving target (STEP only just operational),
language coverage. Distinguish from substantive limitations (kept for the
conclusion chapter).}
